\documentclass[12pt]{article}

\usepackage[a4paper, margin=2cm]{geometry}
\usepackage{fontspec}
\usepackage{polyglossia}
\usepackage{algorithm2e}
\usepackage{setspace}
\usepackage{graphicx}
\usepackage{amsmath}
\usepackage{minted}
\usepackage{float}

\setmainfont{NewCM10-Regular}[
  Path,
  BoldFont        = NewCM10-Bold,
  ItalicFont      = NewCM10-Italic,
  BoldItalicFont  = NewCM10-BoldItalic,
]

\setsansfont{NewCMSans10-Regular}[
  Path,
  BoldFont        = NewCMSans10-Bold,
]

\setmonofont{NewCMMono10-Regular}[
  Path,
  BoldFont        = NewCMMono10-Bold,
  ItalicFont      = NewCMMono10-Italic,
]

\setdefaultlanguage{russian}

\begin{document}
\begin{titlepage}
\centering

{Министерство науки и высшего образования Российской Федерации}
\vskip 2mm
{\large\scshape Новосибирский государственный технический университет}
\vskip 2mm
{Кафедра прикладной математики}

\vfill
\vskip 3cm

{\scshape Лабораторная работа}
\vskip 3mm
\begin{spacing}{1.7}
    {\huge\bf Метод конечных элементов для нелинейных уравнений}
\end{spacing}
\vskip 5mm
{\scriptsize ФПМИ, ПМ-24}
\vskip 2mm
{\large\it Параскун И.}

\vfill

{преподаватели}
\vskip 2mm
{\scshape Вагин Д. В., д.т.н., доцент}
\vskip 1mm
{\scshape Задорожный А. Г., к.т.н., доцент}
\vskip 1mm
{\scshape Леонович Д. А., ассистент}

\vfill

{Новосибирск, \the\year}

\end{titlepage}

\setcounter{page}{2}

\section{Постановка задачи}

\noindent Построить МКЭ для нелинейного уравнения эллиптического типа

\begin{align}
  \label{equation} -\nabla(\lambda(u)\nabla u) + \gamma (u) u = f(u)
\end{align}

\vskip 2mm
\noindent с учётом следующих краевых условий:

\begin{align}
  &u \big|_{S_1} = u_g \\
  &\lambda(u)\frac{\partial u}{\partial n} \big|_{S_2} = \theta(u) \\
  &\lambda(u)\frac{\partial u}{\partial n} \big|_{S_3} +
    \beta(u \big|_{S_3} - u_{\beta}(u)) = 0
\end{align}

\section{Теоретическая часть}

\noindent Эквивалентная вариационная постановка в форме уравнения Галёркина для
уравнения \ref{equation}:

\begin{align}
  \begin{split}
    &\int_{\Omega}\lambda(u)\nabla u \nabla v_0 d\Omega +
    \int_{\Omega}\gamma(u)uv_0 d\Omega +
    \int_{S_3}\beta(u)uv_0 dS \\
    &= \int_{\Omega}f(u)v_0 d\Omega +
    \int_{S_2}\theta(u)v_0 dS +
    \int_{S_3}\beta(u)u_{\beta}(u)v_0 dS, \ \ \forall v_0 \in H_0^1
  \end{split}
\end{align}

\vskip 2mm
\noindent Раскладывая функции $u$ и $v_0$ по базису, переходим 
к конечноэлементной СЛАУ:

\begin{align}
  \begin{split}
    &\sum_{j=1}^{n}\big(
      \int_{\Omega}\lambda(u)\nabla \psi_j \nabla \psi_i d\Omega +
      \int_{\Omega}\gamma(u) \psi_j \psi_i d\Omega +
      \int_{S_3}\beta(u) \psi_j \psi_i dS
    \big)q_j \\
    &= \int_{\Omega}f(u) \psi_i d\Omega +
    \int_{S_2}\theta(u) \psi_i dS +
    \int_{S_3}\beta(u)u_{\beta}(u) \psi_i dS
  \end{split}
\end{align}

\vskip 2mm
\noindent или, в более удобной форме записи

\begin{align*}
  &A_{ij} = G_{ij} + M^{\gamma}_{ij} + M^{S_3}_{ij} \\
  &b_{i} = b^{\Omega}_{i} + b^{S_2}_{i} + b^{S_3}_{i}
\end{align*}

\vskip 2mm
\noindent где

\begin{align}
  &\hat{G}_{ij} = \int_{\hat{\Omega}}\hat{\lambda}(u)\nabla\hat{\psi}_j\nabla\hat{\psi}_i d\hat{\Omega} \\
  &\hat{M}^{\gamma}_{ij} = \int_{\hat{\Omega}}\hat{\gamma}(u)\hat{\psi}_j\hat{\psi}_i d\hat{\Omega} \\
  &\hat{M}^{S_3}_{ij} = \int_{\hat{S_3}}\hat{\beta}(u)\hat{\psi}_j\hat{\psi}_i d\hat{S} \\
  &\hat{b}^{\Omega}_{i} = \int_{\hat{\Omega}}\hat{f}(u)\hat{\psi}_i d\hat{\Omega} \\
  &\hat{b}^{S_2}_{i} = \int_{\hat{S_2}}\hat{\theta}(u)\hat{\psi}_i d\hat{S} \\
  &\hat{b}^{S_3}_{i} = \int_{\hat{S_3}}\hat{\beta}(u)\hat{u}_{\beta}(u)\hat{\psi}_i d\hat{S}
\end{align}

\vskip 2mm
\noindent Раскладывая параметры по базисным функциям на конечном элементе

\begin{align}
  &\hat{G}_{ij} = \int_{\hat{\Omega}}
    \sum_{k=1}^{8}(\hat{\lambda}_{k}\hat{\psi}_k)
    \nabla\hat{\psi}_j
    \nabla\hat{\psi}_i d\hat{\Omega} 
  = \sum_{k=1}^{8}\hat{\lambda}_{k}\int_{\hat{\Omega}}
    \hat{\psi}_k
    \nabla\hat{\psi}_j
    \nabla\hat{\psi}_i d\hat{\Omega} \\
  &\hat{M}^{\gamma}_{ij} = \int_{\hat{\Omega}}
    \sum_{k=1}^{8}(\hat{\gamma}_{k}\hat{\psi}_k)
    \hat{\psi}_j
    \hat{\psi}_i d\hat{\Omega} 
  = \sum_{k=1}^{8}\hat{\gamma}_{k}\int_{\hat{\Omega}}
    \hat{\psi}_k
    \hat{\psi}_j
    \hat{\psi}_i d\hat{\Omega} \\
  &\hat{M}^{S_3}_{ij} = \int_{\hat{S_3}}
    \sum_{k=1}^{4}(\hat{\beta}_{k}\hat{\psi}_k)
    \hat{\psi}_j
    \hat{\psi}_i d\hat{S}
  = \sum_{k=1}^{4}\hat{\beta}_{k}\int_{\hat{S_3}}
    \hat{\psi}_k
    \hat{\psi}_j
    \hat{\psi}_i d\hat{S} \\
  &\hat{b}^{\Omega}_{i} = \int_{\hat{\Omega}}
    \sum_{k=1}^{8}(\hat{f}_{k}\hat{\psi}_k)
    \hat{\psi}_i d\hat{\Omega} 
  = \sum_{k=1}^{8}\hat{f}_{k}\int_{\hat{\Omega}}
    \hat{\psi}_k
    \hat{\psi}_i d\hat{\Omega} \\
  &\hat{b}^{S_2}_{i} = \int_{\hat{S_2}}
    \sum_{k=1}^{4}(\hat{\theta}_{k}\hat{\psi}_k)
    \hat{\psi}_i d\hat{S} 
  = \sum_{k=1}^{4}\hat{\theta}_{k}\int_{\hat{S_2}}
    \hat{\psi}_k
    \hat{\psi}_i d\hat{S} \\
  &\hat{b}^{S_3}_{i} = \int_{\hat{S_3}}
    \sum_{k=1}^{4}(\hat{\beta}_{k}\hat{\psi}_k)
    \sum_{k=1}^{4}(\hat{u}_{\beta, k}\hat{\psi}_k)
    \hat{\psi}_i d\hat{S}
  = \sum_{k=1}^{4}\hat{\beta}_{k}\sum_{j=1}^{4}
    \hat{u}_{\beta, j}
    \int_{\hat{S_3}}\hat{\psi}_k\hat{\psi}_j\hat{\psi}_i d\hat{S}
\end{align}

\vskip 2mm
\noindent Учитывая вид трилинейных базисных функций

\begin{align*}
  & X_1(x) = \frac{x_{p+1}-x}{h_x}, & &X_2(x) = \frac{x - x_{p}}{h_x}, & &h_x = x_{p+1} - x_p, \\
  & Y_1(y) = \frac{y_{s+1}-y}{h_y}, & &Y_2(y) = \frac{y - y_{s}}{h_y}, & &h_y = y_{s+1} - y_s, \\
  & Z_1(z) = \frac{z_{r+1}-z}{h_z}, & &Z_2(z) = \frac{z - z_{r}}{h_z}, & &h_z = z_{r+1} - z_r.
\end{align*}

\begin{align}
  \psi_i = X_{\mu(i)}Y_{\nu(i)}Z_{\vartheta(i)}
\end{align}

\vskip 2mm
\noindent получаем следующие соотношения

\begin{align}
  \begin{split}
    \hat{G}_{ij} = \sum_{k=1}^{8}\hat{\lambda}_{k}(
      &{}^{n}G^{x}_{\mu(k)\mu(j)\mu(i)}{}^{n}M^{y}_{\nu(k)\nu(j)\nu(i)}{}^{n}M^{z}_{\vartheta(k)\vartheta(j)\vartheta(i)}+\\
      &{}^{n}M^{x}_{\mu(k)\mu(j)\mu(i)}{}^{n}G^{y}_{\nu(k)\nu(j)\nu(i)}{}^{n}M^{z}_{\vartheta(k)\vartheta(j)\vartheta(i)}+\\
      &{}^{n}M^{x}_{\mu(k)\mu(j)\mu(i)}{}^{n}M^{y}_{\nu(k)\nu(j)\nu(i)}{}^{n}G^{z}_{\vartheta(k)\vartheta(j)\vartheta(i)}
    )
  \end{split}
\end{align}

\begin{align}
  &\hat{M}^{\gamma}_{ij} = \sum_{k=1}^{8}\hat{\gamma}_{k}(
      {}^{n}M^{x}_{\mu(k)\mu(j)\mu(i)}{}^{n}M^{y}_{\nu(k)\nu(j)\nu(i)}{}^{n}M^{z}_{\vartheta(k)\vartheta(j)\vartheta(i)}
  ) \\
  &\hat{M}^{S_3}_{ij} = \sum_{k=1}^{4}\hat{\beta}_{k}(
      {}^{n}M^{\xi}_{\mu(k)\mu(j)\mu(i)}{}^{n}M^{\chi}_{\nu(k)\nu(j)\nu(i)}
  ) \\
  &\hat{b}^{\Omega}_{i} = \sum_{k=1}^{8}\hat{f}_{k}(
      M^{x}_{\mu(k)\mu(i)}M^{y}_{\nu(k)\nu(i)}M^{z}_{\vartheta(k)\vartheta(i)}
  ) \\
  &\hat{b}^{S_2}_{i} = \sum_{k=1}^{4}\hat{\theta}_{k}(
      M^{\xi}_{\mu(k)\mu(i)}M^{\chi}_{\nu(k)\nu(i)}
  ) \\
  &\hat{b}^{S_3}_{i} = \sum_{k=1}^{4}\hat{\beta}_{k}\sum_{j=1}^{4}\hat{u}_{\beta, j}(
      {}^{n}M^{\xi}_{\mu(k)\mu(j)\mu(i)}{}^{n}M^{\chi}_{\nu(k)\nu(j)\nu(i)}
  )
\end{align}

\section{Описание разработанных программ}

\section{Описание тестирования программ}

\section{Проведенные исследования и выводы}

\section{Тексты основных модулей программ}

\end{document}
